% ============================================================
%  V. EXPERIMENTAL RESULTS
% ============================================================
\section{Experimental Results}
\label{sec:experimental_results}

\subsection{Experimental Setup}

Experiments ran on the physical Spot+Z1 platform with an Orbbec Femto Bolt on
Spot's body and a RealSense D435 at the \ac{EE}, computing across SpotCore
and an external PC over \ac{ROS}~2 Jazzy.  Ten supine patient
positions were defined (\SI{2}{\metre}--\SI{4}{\metre}, $\pm\SI{45}{\degree}$
angular offset); each ran in both modes below, yielding twenty paired trials.

\begin{itemize}
  \item \textbf{Baseline}: sequential stop-and-manipulate. Fixed three-pose scan
        after navigation, no concurrency.
  \item \textbf{\ac{TERESA}}: anticipatory scanner overlapping navigation,
        adaptive grid, and QualityMonitor-modulated velocity.
\end{itemize}

Metrics: total mission time, idle time, scan overlap (fraction of body scan
completed before handoff), \ac{EE} positioning error at each \ac{FAST} point,
and lock time.

\subsection{Search Phase: Whole-Body Coverage and Dual-Sensor Lock}

We isolated the arm's contribution to search through an ablation experiment:
\ac{TERESA} (full system) versus a Spot-only baseline with the Z1 arm disabled
and only the front-facing Orbbec active. Twenty trials were conducted at five
body yaw angles ($0^\circ$--$180^\circ$, two repetitions each), measuring
time to Orbbec posture confidence $\geq$ \num{0.70}.

With the arm's seven scanning poses providing full \SI{360}{\degree} coverage
from a single base position, \ac{TERESA} required no base rotation: frontal
patients were detected near-instantly, the rest during the scan, yielding a
mean of \SI{3.2 \pm 1.1}{\second}. Without the arm, the fixed camera's
\SI{77}{\degree} effective field of view forced five discrete base rotations
($\sim$\SI{7}{\second} each at \SI{0.2}{\radian\per\second}, plus
\SI{2.0}{\second} dwell), giving \SI{14.8 \pm 2.3}{\second}---a
\textbf{5.0}$\times$ slowdown. A \num{10000}-trial Monte Carlo simulation
confirmed this ratio across all patient orientations.

The dual-sensor architecture added robustness independently of speed: in three
of the ten \ac{TERESA} trials, the Orbbec alone failed to lock. These failures
stemmed from the Orbbec's fixed horizontal viewpoint at
\SI{0.5}{\metre} above the ground, which is suboptimal for detecting a supine
body whose visual features---torso, limbs, head---are primarily visible from
above. The arm-mounted RealSense, positioned by the \ac{WBC} look-at controller
to observe the scene from an elevated angle, detected partial torso keypoints
in these cases and guided the base toward the body, giving the Orbbec a second
observation window at a more favourable orientation. Without this redundancy,
these trials would have resumed the coarse search cycle, adding
$\sim$\SI{14}{\second}. The arm thus serves two complementary roles: its
\SI{360}{\degree} coverage eliminates the base-rotation bottleneck, while its
elevated viewpoint provides the sensor redundancy that prevents detection
failures a single-camera system would incur.

\subsection{End-to-End Mission Performance}

Across all twenty end-to-end trials, \ac{TERESA} reduced total mission time
from \SI{100.2 \pm 8.1}{\second} to \SI{89.7 \pm 7.3}{\second} (\SI{10.5}{\percent}).
The gain was concentrated in idle time: in the baseline, the body scan executed
after the base stopped, adding $\sim$\SI{10}{\second} of idle time; \ac{TERESA}
absorbed this into the approach, completing $\sim$\SI{80}{\percent} of the scan
before reaching the handoff point. The adaptive scan grid accelerated this
further: in the majority of trials, torso keypoints were already resolved before
approach, and the grid automatically selected the minimal two-pose pattern
rather than the full three-pose sweep. Pre-planned body reconfiguration reduced
probe positioning error from \SI{44.7 \pm 5.2}{\milli\metre} to
\SI{29.8 \pm 4.1}{\milli\metre}, with the largest improvement at lateral
\ac{FAST} targets. The exposure scan dominated the mission budget at
$\sim$\SI{40}{\second} in both configurations, limiting the overall speedup.
In manual operation, operators could skip unreachable grid points, reducing
total time to $\sim$\SI{70}{\second} at the cost of incomplete coverage.

\subsection{Whole-Body Reconfiguration}

The body posture optimizer described in Section~\ref{sec:body_reconf} proved
essential for reaching distant anatomical targets on a full-scale body.  We
compared four configurations of increasing capability.  The fixed nominal
posture (height \SI{0}{\metre}, pitch \SI{0}{\degree}, no base displacement),
representing Spot's default standing pose, reached only
$\sim$\SI{40}{\percent} of exposure grid points and $\sim$\SI{50}{\percent}
of \ac{FAST} targets.  A heuristic posture---moderately lowered
($\sim$\SI{-0.15}{\metre}) with a slight forward pitch
($\sim$\SI{5}{\degree})---improved these to $\sim$\SI{55}{\percent} and
$\sim$\SI{65}{\percent}, but could not adapt to individual point locations.
The 2D optimiser, selecting the best $(h, \varphi)$ per point via grid search
with \ac{IK}-driven retry, raised completion to $\sim$\SI{70}{\percent}
(exposure) and $\sim$\SI{85}{\percent} (\ac{FAST}).  Finally, the 3D
optimiser added longitudinal base walking $dy$ with the regularised cost
$J = \|\mathbf{p} - \mathbf{p}_{\mathrm{sweet}}\| + \lambda\,|dy|$
($\lambda = 0.50$), achieving $\sim$\SI{85}{\percent} and
\SI{100}{\percent} respectively.  The penalty encouraged height and pitch
adjustments over walking---reducing energy consumption and producing smoother
motion---while still allowing displacement when necessary.  The residual
\SI{15}{\percent} of exposure points were at the furthest extremities (head,
feet), beyond the arm's maximum reach even with full base repositioning;
these were skipped gracefully.  Each tier contributed measurably to mission
completeness, with the movement penalty ensuring that base walking is used
only when strictly needed.

\subsection{Exposure Body Scanning and Injury Detection}

The exposure pipeline was evaluated with XX trials on supine patients,
generating XX~$\pm$~XX scan points. Per-point dwell of
\SI{2}{\second} yielded scan durations of XX~$\pm$~XX\,s. The injury detection pipeline achieved a wound
recall of XX\,\% and burn classification accuracy of XX\,\%,
with a false positive rate of XX per scan on average. The 3D localisation error
of detected injuries was XX~$\pm$~XX\,mm.

\subsection{Discussion}

The results demonstrate that treating body posture as an active perceptual
resource yields measurable improvements in both mission speed and
completeness.  The arm's \SI{360}{\degree} coverage eliminates the
base-rotation bottleneck during search, while body reconfiguration lifts
completion from $\sim$\SI{40}{\percent} to near-\SI{100}{\percent}.  The
dual-sensor architecture prevents detection failures that a single-camera
system would incur.  Several limitations remain: the system assumes a
single stationary patient on flat terrain; ultrasound image quality and
clinical scanning pressure were not evaluated; and the injury detection
pipeline relies on pre-trained models whose performance may vary across
wound types and skin tones.  These aspects are the subject of ongoing work.
